\section{Competências Desenvolvidas pelo Egresso}

O \NOMECURSO fundamenta sua formação em um modelo pedagógico orientado por competências, o qual articula conhecimentos científicos, habilidades técnicas e atitudes ético-profissionais. Essa abordagem visa garantir a formação de engenheiros capazes de atuar com autonomia intelectual, responsabilidade social e domínio tecnológico em ambientes industriais complexos e dinâmicos.  

As competências são estruturadas segundo as Diretrizes Curriculares Nacionais de Engenharia (Resolução CNE/CES nº 2/2019), assegurando que o egresso desenvolva visão sistêmica, capacidade de inovação e atuação colaborativa. Nessas, o aprendizado ocorre de modo integrado entre atividades teóricas, práticas e extensionistas, consolidando a articulação entre o saber, o saber fazer e o saber ser.  

%--------------------------------------------------------------
Os conhecimentos, habilidades e atitudes necessários ao desenvolvimento das
competências foram elaborados a partir do percurso de aprendizagem necessário à construção de saberes essenciais a cada elemento de competência, considerando sua mensuração, as experiências de aprendizagem e os mecanismos de avaliação.

As perguntas-guia que assistiram esse processo, baseadas em \cite{res259}, foram:

\begin{enumerate}
		\item Quais estratégias didáticas e recursos conceituais devem ser mobilizados em sala de aula para favorecer a compreensão profunda de determinado conceito pelo estudante?
		
		\item Considerando que cada estudante possui repertórios prévios e modos singulares de aprendizagem, de que forma o planejamento pedagógico pode estruturar experiências formativas que respeitem e integrem essas particularidades?
		
		\item Quais tipos de experiências concretas, contextualizadas e sequenciadas são necessárias para promover o desenvolvimento efetivo de uma habilidade específica no estudante?
		
		\item Quais ações, procedimentos ou rotinas acadêmicas o estudante deve realizar ao iniciar a participação em um encontro presencial, de modo a potencializar seu engajamento cognitivo?
		
		\item Quais métodos e instrumentos avaliativos são adequados para mensurar, de maneira objetiva e sistemática, os resultados oriundos das experiências de aprendizagem propostas?
		
		\item Por meio de quais indicadores e critérios avaliativos é possível monitorar, ao longo do tempo, o progresso do estudante no desenvolvimento de uma habilidade específica?
	\end{enumerate}
%	--------------------------------------------------------------

Isto posto, na Tab. \ref{tab:competencias_gerais}  é apresentado o conjunto de competências gerais que constitui a base formativa do egresso e orienta a construção das competências específicas do curso.


\begin{landscape}
	
	\scriptsize
	\begin{longtable}{m{1cm} m{5cm} m{8cm} m{5cm} m{5cm}}
		\caption{Competências gerais desenvolvidas no curso e seus desdobramentos}
		\label{tab:competencias_gerais} \\
		\hline
		\textbf{Cód} & \textbf{Competências Gerais} & \textbf{Objetivos de Aprendizagem / Habilidades} & \textbf{Conhecimento} & \textbf{Atitude} \\
		\hline
		\endfirsthead
		
		\hline
		\textbf{Cód} & \textbf{Competências Gerais} & \textbf{Objetivos de Aprendizagem / Habilidades} & \textbf{Conhecimento} & \textbf{Atitude} \\
		\hline
		\endhead
		
		\multicolumn{5}{r}{\textit{Continua na próxima página}}\\
		\endfoot
		
		\hline
		\caption*{\\ Fonte: \normalfont Adaptado de \citeonline{res259}}
		\endlastfoot
		
		
		% --------------------- CG1 ----------------------
		\multirow{2}{*}{CG1} &
		\multirow{2}{5cm}{Formular e conceber soluções desejáveis de engenharia, analisando e compreendendo os usuários dessas soluções e seu contexto.}
		& HG11 – Ser capaz de utilizar técnicas adequadas de observação, compreensão, registro e análise das necessidades dos usuários e de seus contextos sociais, culturais, legais, ambientais e econômicos.
		& Conteúdo deste conhecimento
		& Conteúdo teste de atitude CE1
		\\ \cline{3-5}
		
		&
		& HG12 – Formular, de maneira ampla e sistêmica, questões de engenharia considerando o usuário e seu contexto, concebendo soluções criativas com técnicas adequadas.
		& Conteúdo deste conhecimento
		& Conteúdo teste de atitude CE1
		\\ \hline
		
		
		% --------------------- CG2 ----------------------
		\multirow{4}{*}{CG2} &
		\multirow{4}{5cm}{Analisar e compreender os fenômenos físicos e químicos por meio de modelos simbólicos, físicos e outros, verificados e validados por experimentação.}
		& HG21 – Ser capaz de modelar fenômenos, sistemas físicos e químicos utilizando ferramentas matemáticas, estatísticas, computacionais e de simulação.
		& 
		& 
		\\ \cline{3-5}
		
		&
		& HG22 – Prever os resultados dos sistemas por meio dos modelos.
		&
		&
		\\ \cline{3-5}
		
		&
		& HG23 – Conceber experimentos que gerem resultados reais dos fenômenos e sistemas estudados.
		&
		&
		\\ \cline{3-5}
		
		&
		& HG24 – Verificar e validar modelos por meio de técnicas adequadas.
		&
		&
		\\ \hline
		
		
		% --------------------- CG3 ----------------------
		\multirow{3}{*}{CG3} &
		\multirow{3}{5cm}{Conceber, projetar e analisar sistemas, produtos, componentes ou processos.}
		& HG31 – Ser capaz de conceber e projetar soluções criativas, desejáveis e viáveis técnica e economicamente.
		&
		&
		\\ \cline{3-5}
		
		&
		& HG32 – Projetar e determinar parâmetros construtivos e operacionais das soluções de engenharia.
		&
		&
		\\ \cline{3-5}
		
		&
		& HG33 – Aplicar conceitos de gestão para planejar, supervisionar, elaborar e coordenar projetos e serviços de Engenharia.
		&
		&
		\\ \hline
		
		
		% --------------------- CG4 ----------------------
		\multirow{5}{*}{CG4} &
		\multirow{5}{5cm}{Implantar, supervisionar e controlar soluções de engenharia.}
		& HG41 – Aplicar conceitos de gestão para planejar, supervisionar e coordenar a implantação das soluções.
		&
		&
		\\ \cline{3-5}
		
		&
		& HG42 – Gerir força de trabalho e recursos físicos, materiais e de informação.
		&
		&
		\\ \cline{3-5}
		
		&
		& HG43 – Desenvolver sensibilidade global nas organizações.
		&
		&
		\\ \cline{3-5}
		
		&
		& HG44 – Projetar e desenvolver estruturas empreendedoras e soluções inovadoras.
		&
		&
		\\ \cline{3-5}
		
		&
		& HG45 – Realizar avaliação crítico-reflexiva dos impactos das soluções de engenharia nos contextos social, legal, econômico e ambiental.
		&
		&
		\\ \hline
		
		
		% --------------------- CG5 ----------------------
		\multirow{1}{*}{CG5} &
		\multirow{1}{5cm}{Comunicar-se eficazmente nas formas escrita, oral e gráfica.}
		& HG51 – Ser capaz de expressar-se adequadamente em língua pátria ou estrangeira, inclusive com TDICs, mantendo-se atualizado em métodos e tecnologias.
		&
		&
		\\ \hline
		
		
		% --------------------- CG6 ----------------------
		\multirow{5}{*}{CG6} &
		\multirow{5}{5cm}{Trabalhar e liderar equipes multidisciplinares.}
		& HG61 – Interagir com diferentes culturas em equipes presenciais ou a distância.
		&
		&
		\\ \cline{3-5}
		
		&
		& HG62 – Atuar de forma colaborativa, ética e profissional em equipes multidisciplinares locais e em rede.
		&
		&
		\\ \cline{3-5}
		
		&
		& HG63 – Gerenciar projetos e liderar proativamente, definindo estratégias e construindo consensos.
		&
		&
		\\ \cline{3-5}
		
		&
		& HG64 – Reconhecer e conviver com diferenças socioculturais nos diversos contextos.
		&
		&
		\\ \cline{3-5}
		
		&
		& HG65 – Preparar-se para liderar empreendimentos nos aspectos de produção, finanças, pessoal e mercado.
		&
		&
		\\ \hline
		
		
		% --------------------- CG7 ----------------------
		\multirow{2}{*}{CG7} &
		\multirow{2}{5cm}{Conhecer e aplicar com ética a legislação e normas da profissão.}
		& HG71 – Compreender legislação, ética e responsabilidade profissional, avaliando impactos da engenharia na sociedade e no ambiente.
		&
		&
		\\ \cline{3-5}
		
		&
		& HG72 – Atuar sempre com ética e respeito à legislação, zelando por sua aplicação no contexto profissional.
		&
		&
		\\ \hline
		
		
		% --------------------- CG8 ----------------------
		\multirow{2}{*}{CG8} &
		\multirow{2}{5cm}{Aprender de forma autônoma e lidar com situações complexas, atualizando-se com avanços da ciência, tecnologia e inovação.}
		& HG81 – Assumir atitude investigativa e autônoma para aprendizagem contínua, produção de novos conhecimentos e desenvolvimento tecnológico.
		&
		&
		\\ \cline{3-5}
		
		&
		& HG82 – Aprender a aprender.
		&
		&
		\\ \hline
		
	\end{longtable}
	
\end{landscape}

%\begin{landscape}
%	
%	\scriptsize
%	\setlength{\tabcolsep}{3pt}
%	
%	\begin{longtable}{m{1cm} m{4.8cm} m{11cm} m{4.2cm} m{4.2cm}}
%		\caption{Competências gerais desenvolvidas no curso e seus desdobramentos}
%		\label{tab:competencias_gerais} \\
%		\hline
%		\textbf{Cód} & \textbf{Competências Gerais} & \textbf{Objetivos de Aprendizagem / Habilidades} & \textbf{Conhecimento} & \textbf{Atitude} \\
%		\hline
%		\endfirsthead
%		
%		\hline
%		\textbf{Cód} & \textbf{Competências Gerais} & \textbf{Objetivos de Aprendizagem / Habilidades} & \textbf{Conhecimento} & \textbf{Atitude} \\
%		\hline
%		\endhead
%		
%%		\hline
%		\multicolumn{5}{r}{\textit{Continua na próxima página}}\\
%		\endfoot
%		
%		\hline
%		\caption*{\\ Fonte: \normalfont Adaptado de \citeonline{res259}}
%		\endlastfoot
%		
%		% --------------------- CG1 ----------------------
%		\multirow{1}{*}{CG1} &
%		Formular e concebersoluções desejáveis de engenharia, analisando e compreendendo os usuários dessas soluções e seu contexto.
%		 &
%		\begin{itemize}
%			\item HG11 – Ser capaz de utilizar técnicas adequadas de observação, compreensão, registro e análise das necessidades dos
%		usuários e de seus contextos sociais, culturais, legais, ambientais e econômicos; 
%		\item HG12 - Formular, de maneira ampla e sistêmica, questões de engenharia, considerando o usuário e seu contexto, concebendo soluções criativas, bem como o uso de técnicas adequadas;
%				\end{itemize}
%				&
%		\multirow{1}{5.2cm}{\begin{itemize}
%				\item Conteúdo deste conhecimento
%				\item Conteúdo deste conhecimento
%		\end{itemize}} &
%		
%		\multirow{1}{5.2cm}{\begin{itemize}
%				\item Conteúdo teste de atitude CE1
%		\end{itemize}} \\
%				\hline
%				
%		% --------------------- CG2 ----------------------
%		\multirow{5}{*}{CG2} &
%		Analisar e compreender os fenômenos físicos e químicos por meio de modelos simbólicos, físicos e outros, verificados e validados por experimentação &
%		HG 21 - Ser capaz de modelar os fenômenos, os sistemas físicos e químicos, utilizando as ferramentas matemáticas,
%		estatísticas, computacionais e de simulação, entre outras.\\
%		&&
%		HG22 - Prever os resultados dos sistemas por meio dos modelos;\\
%		&&
%		HG23 - Conceber experimentos que gerem resultados reais para o comportamento dos fenômenos e sistemas em estudo.\\
%		&&
%		HG24 - Verificar e validar os modelos por meio de técnicas adequadas;\\
%		\hline
%		
%		% --------------------- CG3 ----------------------
%		\multirow{4}{*}{CG3} &
%		Conceber, projetar e analisar sistemas, produtos (bens e serviços), componentes ou processos. &
%		HG31 - Ser capaz de conceber e projetar soluções criativas, desejáveis e viáveis, técnica e economicamente, nos contextos em que serão aplicadas.\\
%		&&
%		HG32 - Projetar e determinar os parâmetros construtivos e operacionais para as soluções de Engenharia.\\
%		&&
%		HG33 - Aplicar conceitos de gestão para planejar, supervisionar, elaborar e coordenar projetos e serviços de
%		Engenharia.\\
%		\hline
%		
%		% --------------------- CG4 ----------------------
%		\multirow{5}{*}{CG4} &
%		Implantar, supervisionar e controlar soluções de engenharia. &
%		HG41 - Ser capaz de aplicar os conceitos de gestão para planejar, supervisionar, elaborar e coordenar a implantação das soluções de Engenharia.\\
%		&&
%		HG42 - Estar apto a gerir, tanto a força de trabalho quanto os recursos físicos, no que diz respeito aos materiais e à informação;\\
%		&&
%		HG43 - Desenvolver sensibilidade global nas organizações;\\
%		&&
%		HG44 - Projetar e desenvolver novas estruturas empreendedoras e soluções inovadoras para os problemas.\\
%		&&
%		HG45 - Realizar a avaliação crítico-reflexiva dos impactos das soluções de Engenharia nos contextos social, legal, econômico e ambiental.\\
%		\hline
%		
%		% --------------------- CG5 ----------------------
%		CG5 &
%		Comunicar-se eficazmente nas formas escrita, oral e gráfica. &
%		HG51 - Ser capaz de expressar-se adequadamente, seja na língua pátria ou em idioma diferente do Português,
%		inclusive por meio do uso consistente das tecnologias digitais de informação e comunicação (TDICs), mantendo-se sempre atualizado em termos de métodos e tecnologias disponíveis. \\
%		\hline
%		
%		% --------------------- CG6 ----------------------
%		\multirow{5}{*}{CG6} &
%		Trabalhar e liderar equipes multidisciplinares. &
%		HG61 - Ser capaz de interagir com as diferentes culturas, mediante o trabalho em equipes presenciais ou a distância,
%		de modo que facilite a construção coletiva.\\
%		&&
%		HG62 - Atuar, de forma colaborativa, ética e profissional em equipes multidisciplinares, tanto localmente quanto em rede;\\
%		&&
%		HG63 - Gerenciar projetos e liderar, de forma proativa e colaborativa, definindo as estratégias e construindo o
%		consenso nos grupos;\\
%		&&
%		HG64 - Reconhecer e conviver com as diferenças socioculturais nos mais diversos níveis em todos os contextos em que atua (globais/locais);\\
%		&&
%		HG65 - Preparar-se para liderar empreendimentos em todos os seus aspectos de produção, de finanças, de pessoal e de mercado;\\
%		\hline
%		
%		% --------------------- CG7 ----------------------
%		\multirow{2}{*}{CG7} &
%		Conhecer e aplicar com ética a legislação e normas da profissão. &
%		HG71 - er capaz de compreender a legislação, a ética e a responsabilidade profissional e avaliar os impactos das
%		atividades de Engenharia na sociedade e no meio ambiente.\\
%		&&
%		HG72 - Atuar sempre respeitando a legislação, e com ética em todas as atividades, zelando para que isto ocorra
%		também no contexto em que estiver atuando; \\
%		\hline
%		
%		% --------------------- CG8 ----------------------
%		\multirow{2}{*}{CG8} &
%		Aprender de forma autônoma e lidar com situações e contextos complexos, atualizando-se em relação aos avanços da ciência, da tecnologia e aos desafios da inovação &
%		HG81 - Ser capaz de assumir atitude investigativa e autônoma, com vistas à aprendizagem contínua, à produção de novos conhecimentos e ao desenvolvimento de novas tecnologias.\\
%		&&	
%		HG82 - Aprender a aprender.\\
%		\hline
%		
%	\end{longtable}
%	
%\end{landscape}



De modo símile, na Tab. \ref{tab:competencias_especificas}  é apresentado o conjunto das competências específicas desenvolvidas ao longo do curso.
\begin{landscape}
	
	\scriptsize
	\setlength{\tabcolsep}{4pt}
	
	\begin{longtable}{m{1.2cm} m{5cm} m{8.2cm} m{5.2cm} m{5.2cm}}
		\caption{Competências específicas desenvolvidas no curso}
		\label{tab:competencias_especificas} \\
		\hline
		\textbf{Cód} & \textbf{Competências Específicas} & \textbf{Objetivos de Aprendizagem / Habilidades} & \textbf{Conhecimento} & \textbf{Atitude} \\
		\hline
		\endfirsthead
		
		\hline
		\textbf{Cód} & \textbf{Competências Específicas} & \textbf{Objetivos de Aprendizagem / Habilidades} & \textbf{Conhecimento} & \textbf{Atitude} \\
		\hline
		\endhead
		
		\multicolumn{5}{r}{\textit{Continua na próxima página}} \\
		\endfoot
		
		\hline
		\caption*{\\Fonte: \normalfont Adaptado de \citeonline{res259}}
		\endlastfoot
		
		
		%%%%%%%%%%%%%%%%%%%%%%%%%%%%%%%% CE1 %%%%%%%%%%%%%%%%%%%%%%%%%%%%%%%%
		
		\multirow{2}{*}{CE1} &
		\multirow{2}{5cm}{Analisar, compreender, selecionar e sugerir materiais metálicos, não-metálicos e compósitos} &
		
		\begin{itemize}
			\item HE11 – Identificar materiais para construção mecânica e compreender sua aplicação em projetos.
			\item HE12 – Realizar ensaios com materiais, relacionando microestrutura e propriedades.
		\end{itemize}
		
		& 
		\multirow{2}{5.2cm}{\begin{itemize}
				\item Conteúdo teste de conhecimento CE1
		\end{itemize}} &
		
		\multirow{2}{5.2cm}{\begin{itemize}
				\item Conteúdo teste de atitude CE1
		\end{itemize}} \\
		
		\hline
		
		
		%%%%%%%%%%%%%%%%%%%%%%%%%%%%%%%% CG2 %%%%%%%%%%%%%%%%%%%%%%%%%%%%%%%%
		
		\multirow{1}{*}{CE2} &
		\multirow{1}{5cm}{-----------} &
		
		\begin{itemize}
			\item HE21 –  ---
			\item HE22 – ---.
		\end{itemize}
		
		&
		\multirow{1}{5.2cm}{\begin{itemize}
				\item --- 
		\end{itemize}} &
		
		\multirow{1}{5.2cm}{\begin{itemize}
				\item --- 
		\end{itemize}} \\
		
	
	\end{longtable}
	
\end{landscape}


